\section{Introducción}
El manejo de incidentes supone un reto importante para las instituciones, sean públicas o 
privadas. En muchos casos, los trámites se llevan a cabo manualmente, mediante correos 
electrónicos dispersos, llamadas telefónicas o mediante whatsapp. Esto conduce a una 
falta de trazabilidad, pérdida de información y retrasos importantes en la atención a las 
solicitudes. Esta situación, expuesta en el Informe de Especificación de Requerimientos 
del Sistema de Gestión de Tickets [10],sirvió como punto de partida para la propuesta de 
un proyecto con el propósito de digitalizar y mejorar estos procedimientos. 

El propósito del sistema propuesto es atender los siguientes requerimientos: 
centralización de datos, alertas en tiempo real, paneles con métricas confiables y la 
posibilidad de personalizar según los roles de los usuarios (técnicos, administradores, 
superadministradores y funcionarios). La arquitectura tecnológica establecida, que 
combina un sitio web basado en React con un backend construido en Laravel y una 
aplicación móvil para técnicos, permite garantizar una comunicación ininterrumpida y un 
monitoreo riguroso de las solicitudes. 

En el marco teórico, se consultaron artículos académicos que sirven como referencia para 
justificar las decisiones tecnológicas y metodológicas. Por ejemplo, el estudio sobre 
Laravel [1] recalca la relevancia" o "en el estudio sobre Laravel [1] se recalca la 
relevancia de este framework en términos de seguridad y rapidez, mientras que el trabajo 
comparativo entre Laravel y Symfony [7] permitió confirmar que, aunque Symfony 
ofrece mayor robustez empresarial, Laravel resulta más accesible y flexible para un 
equipo en formación. Del mismo modo, React se presenta como una biblioteca flexible y 
fácil de adoptar, lo que coincide con investigaciones que la destacan como la opción 
preferida frente a Angular y Vue en numerosos proyectos [2], [4], [5]. React Native, por 
su parte, constituye una herramienta estratégica para futuras ampliaciones hacia entornos 
multiplataforma [3]. 

El proyecto se nutrió también de los aportes de metodologías ágiles. El marco de trabajo 
Scrum se utilizó como base para organizar el flujo de trabajo, asegurando entregas 
incrementales y revisiones continuas [6]. Sin embargo, tal como lo advierte la literatura 
sobre metodologías ágiles [9], la adopción de estos enfoques debe complementarse con 
unas buenas prácticas de ingeniería de software para evitar problemas y asegurar la 
escalabilidad del sistema. 

El objetivo de este artículo es no solo presentar los progresos logrados en el desarrollo 
del sistema de gestión de tickets, sino también llevar a cabo un análisis comparativo con 
la literatura disponible para confirmar las decisiones tomadas en términos tecnológicos y 
metodológicos y establecer áreas para futuras investigaciones. 